\documentclass{notaaula}

\titulonota{Corrente elétrica alternada}
\creditos{3º ano · Física · Agosto/2026 · Prof. Josué Cavalcante}

\begin{document}
\cabecalho

\begin{sobre}
Eletrodinâmica · noções de CA · $i(t)$ com seno e cosseno · 3º A/B/C · FGB.
\end{sobre}

% ============================================================
\section{Corrente contínua e corrente alternada}

\begin{copiar}{Noção de CA}
\definicao Na corrente \dest{contínua} (CC), a intensidade não troca de sinal
com o tempo: as cargas seguem \dest{sempre o mesmo sentido} no fio. Na corrente
\dest{alternada} (CA), a intensidade \dest{muda de sinal} periodicamente: o
sentido no fio \dest{inverte} a cada meio ciclo. Nesta aula descrevemos só
\textit{como} $i$ varia no tempo, com seno e cosseno.
\diaadia{A tomada da casa entrega CA. A pilha, no circuito simples, entrega CC.}
\end{copiar}

\begin{figuranota}{Eixo horizontal: tempo $t$. Mesma amplitude $i_0$: em azul,
CC constante; em vermelho, CA que oscila e troca de sinal.}
\begin{tikzpicture}
\begin{axis}[
  width=1.06\linewidth, height=4.3cm,
  axis lines=middle, axis line style={naSuave, -{Stealth[length=4pt]}},
  xmin=-0.12, xmax=2.22, ymin=-1.55, ymax=1.95,
  xtick={0,0.5,1,1.5,2}, xticklabels={$0$,$T/2$,$T$,$3T/2$,$2T$},
  ytick={-1,1}, yticklabels={$-i_0$,$i_0$},
  tick style={naSuave},
  tick label style={font=\tiny, color=naSuave},
  xlabel={$t$ (tempo)}, xlabel style={font=\tiny, color=naSuave},
  clip=false,
]
  \addplot[naAzul,  very thick, domain=0:2, samples=2]   {1};
  \addplot[naCoral, very thick, domain=0:2, samples=300] {sin(deg(2*pi*x))};
  \node[font=\tiny, color=naAzul,  anchor=west] at (axis cs:0.05,1.28) {CC\ \ $i=i_0$};
  \node[font=\tiny, color=naCoral, anchor=west] at (axis cs:0.95,1.66) {CA\ \ $i=i_0\sen(\omega t)$};
\end{axis}
\end{tikzpicture}
\end{figuranota}

% ============================================================
\section{Radiano e conversão para graus}

\begin{copiar}{Radiano}
\definicao O \dest{radiano} mede ângulo pelo arco: se um arco tem comprimento
$s$ sobre um círculo de raio $R$,
\[ \theta = \frac{s}{R} \quad \text{(em rad)} \]
Quando $s=R$, $\theta = 1$ rad. A volta inteira tem arco $2\pi R$, logo
\[ 2\pi\ \mathrm{rad} = 360^\circ \qquad \pi\ \mathrm{rad} = 180^\circ \]
\simbolos{$\theta$ (rad ou °); $s$ (m); $R$ (m).}
\end{copiar}

\begin{copiar}{Graus $\leftrightarrow$ radianos}
\formulas
\[ \theta^\circ = \theta_{\mathrm{rad}}\cdot\frac{180}{\pi}
   \qquad
   \theta_{\mathrm{rad}} = \theta^\circ\cdot\frac{\pi}{180} \]
Valores que você vai usar o tempo todo:
\[ 30^\circ = \frac{\pi}{6},\quad 45^\circ = \frac{\pi}{4},\quad
   60^\circ = \frac{\pi}{3},\quad 90^\circ = \frac{\pi}{2}. \]
\atencao{Na fórmula $i = i_0\sen(\omega t)$, o produto $\omega t$ entra em
\dest{radianos}. Se o calculador estiver em graus, o seno sai errado.}
\end{copiar}

\begin{figuranota}{Ângulo de 1 rad: o arco vermelho tem o mesmo comprimento do
raio. À direita, a conversão.}
\begin{tikzpicture}[scale=0.86]
  \def\R{1.75}
  \draw[naTexto, thick] (0,0) circle (\R);
  \draw[naAzul, thick] (0,0) -- (0:\R);
  \draw[naAzul, thick] (0,0) -- (57.2958:\R);
  \draw[naCoral, very thick] (0:\R) arc (0:57.2958:\R);
  \node[font=\tiny, color=naSuave] at (0:{\R/2}) [below] {$R$};
  \node[font=\tiny, color=naCoral] at (28.6:{\R+0.30}) {$s=R$};
  \node[font=\tiny, color=naSuave] at (28.6:0.62) {rad};
  % legenda à direita
  \begin{scope}[shift={(2.55,1.35)}, font=\tiny]
    \node[anchor=west, color=naTexto, font=\tiny\bfseries] at (0,0.55) {volta completa};
    \node[anchor=west, color=naVerde] at (0,0.05) {$2\pi$ rad $=360^\circ$};
    \node[anchor=west, color=naVerde] at (0,-0.35) {$\pi$ rad $=180^\circ$};
    \node[anchor=west, color=naSuave] at (0,-1.05) {graus $=$ rad $\cdot\,180/\pi$};
    \node[anchor=west, color=naSuave] at (0,-1.45) {rad $=$ graus $\cdot\,\pi/180$};
    \node[anchor=west, color=naAmbar] at (0,-2.15) {$1$ rad $\approx 57{,}3^\circ$};
    \node[anchor=west, color=naAmbar] at (0,-2.55) {$1^\circ \approx 0{,}0175$ rad};
  \end{scope}
\end{tikzpicture}
\end{figuranota}

% ============================================================
\section{Funções seno e cosseno}

\begin{copiar}{Seno e cosseno}
\rotulo{No triângulo retângulo} (ângulo agudo $\alpha$):
\[
\begin{aligned}
  \sen\alpha &= \frac{\text{cateto oposto}}{\text{hipotenusa}},\\[2pt]
  \cos\alpha &= \frac{\text{cateto adjacente}}{\text{hipotenusa}}.
\end{aligned}
\]
No círculo de raio $1$, $\sen\theta$ é a ordenada e $\cos\theta$ é a abscissa
do ponto no arco $\theta$.

\rotulo{Ciclo e sinal.} Um período é $2\pi$ rad ($360^\circ$).
\begin{center}\footnotesize
\begin{tabular}{r|cccc}
\toprule
$\theta$      & $0$  & $\pi/2$ ($90^\circ$) & $\pi$ ($180^\circ$) & $3\pi/2$ ($270^\circ$)\\
\midrule
$\sen\theta$  & $0$  & $+1$ & $0$  & $-1$\\
$\cos\theta$  & $+1$ & $0$  & $-1$ & $0$\\
\bottomrule
\end{tabular}
\end{center}
O cosseno está adiantado de $\pi/2$ em relação ao seno:
$\cos\theta = \sen(\theta + \pi/2)$.
\atencao{Seno (ou cosseno) \dest{negativo} não é erro: na CA, $i<0$ é o sentido
contrário ao positivo combinado.}
\end{copiar}

\begin{figuranota}{Um período do seno: zeros em $0$, $\pi$ e $2\pi$; máximo em
$\pi/2$; mínimo em $3\pi/2$.}
\begin{tikzpicture}
\begin{axis}[
  width=1.06\linewidth, height=4.3cm,
  axis lines=middle, axis line style={naSuave, -{Stealth[length=4pt]}},
  xmin=-0.35, xmax=7.1, ymin=-1.70, ymax=1.6,
  xtick=\empty, ytick={-1,1}, yticklabels={$-1$,$+1$},
  tick style={naSuave}, tick label style={font=\tiny, color=naSuave},
  xlabel={$\omega t$}, xlabel style={font=\tiny, color=naSuave},
  clip=false,
]
  \addplot[naVerde, very thick, domain=0:2*pi, samples=300] {sin(deg(x))};
  \foreach \x/\lab in {0/$0$, 1.5708/$\pi/2$, 3.1416/$\pi$,
                       4.7124/$3\pi/2$, 6.2832/$2\pi$} {
    \edef\temp{\noexpand\addplot[naCoral, mark=*, mark size=1.3pt, only marks]
      coordinates {(\x, {sin(deg(\x))})};}\temp
  }
  \node[font=\tiny, color=naSuave] at (axis cs:0,-0.30)      {$0$};
  \node[font=\tiny, color=naSuave] at (axis cs:1.5708,1.28)  {$\pi/2$};
  \node[font=\tiny, color=naSuave] at (axis cs:3.1416,-0.30) {$\pi$};
  \node[font=\tiny, color=naSuave] at (axis cs:4.7124,-0.66) {$3\pi/2$};
  \node[font=\tiny, color=naSuave] at (axis cs:6.2832,0.28)  {$2\pi$};
  \draw[naAzul, {Stealth[length=4pt]}-{Stealth[length=4pt]}]
    (axis cs:0,-1.52) -- (axis cs:6.2832,-1.52)
    node[midway, below, font=\tiny, color=naAzul] {$T$};
\end{axis}
\end{tikzpicture}
\end{figuranota}

% ============================================================
\section{Intensidade instantânea \texorpdfstring{$i(t)$}{i(t)}}

\begin{copiar}{Corrente senoidal}
\modelo A intensidade \dest{instantânea} de uma CA senoidal é
\[ i(t) = i_0\sen(\omega t) \]
$i_0 > 0$ é a \dest{amplitude} (valor máximo de $|i|$). O argumento $\omega t$
está em radianos.

Se a oscilação começa no máximo ($t=0 \Rightarrow i=i_0$), usamos o cosseno:
\[ i(t) = i_0\cos(\omega t) \]
As duas descrevem a \dest{mesma} forma de onda, só mudam o instante em que
começamos a contar o tempo.
\simbolos{$i$ (A); $i_0$ (A); $\omega$ (rad/s); $t$ (s).}
\end{copiar}

\begin{copiar}{Pulso, frequência e período}
\relacoes
\[ \omega = 2\pi f = \frac{2\pi}{T} \]
Logo as formas equivalentes
\[ i(t) = i_0\sen(2\pi f t) \qquad i(t) = i_0\sen\!\left(\frac{2\pi}{T}\,t\right) \]
$T$ é o \dest{período} (tempo de um ciclo); $f = 1/T$ é a \dest{frequência}
(ciclos por segundo, hertz).
\diaadia{Na rede residencial brasileira, $f = 60\un{Hz}$, então
$T = 1/60\un{s}$.}
\end{copiar}

% ============================================================
\section{Instantes especiais}

Com $i(t) = i_0\sen(\omega t)$:
\begin{itens}
\item $\omega t = 0$ ou $\pi$ ou $2\pi$: $\sen = 0 \Rightarrow i = 0$
      (passa pelo zero).
\item $\omega t = \pi/2$: $\sen = 1 \Rightarrow i = +i_0$
      (pico no sentido positivo).
\item $\omega t = 3\pi/2$: $\sen = -1 \Rightarrow i = -i_0$
      (pico no sentido contrário).
\end{itens}

\begin{exemplo}
Uma CA tem $i_0 = \dec{4,0}\un{A}$ e $f = 60\un{Hz}$, com
$i(t) = i_0\sen(2\pi f t)$.
\begin{alternativas}
\item Escreva $i(t)$ com números.
\item Ache $i$ em $t=0$ e em $t = T/4$.
\item Ache o período $T$.
\end{alternativas}
\resolucao (a) $2\pi f = 2\pi\cdot 60 = 120\pi$, logo
\resultado{$i(t) = \dec{4,0}\sen(120\pi\,t)$} (com $t$ em segundos).

(b) Em $t=0$, $\sen 0 = 0 \Rightarrow i = 0$. $T = 1/60\un{s}$, então
$t = T/4$ corresponde a $\omega t = 2\pi f\cdot T/4 = 2\pi/4 = \pi/2$,
$\sen(\pi/2) = 1 \Rightarrow i = +\dec{4,0}\un{A}$.

(c) \resultado{$T = 1/f = 1/60\un{s} \approx \dec{0,017}\un{s}$.}
\end{exemplo}

\begin{dica}
Não troque $i_0$ (pico) pela leitura “do dia a dia” da tomada: o número
$127\un{V}$ da rede é um \dest{valor eficaz}, menor que o pico. Nesta aula o
foco é $i(t)$ e o seno — o valor eficaz volta com calma.
\end{dica}

% ============================================================
\begin{exercicios}
\nivel{1}{Conceitos}
\begin{questoes}
\item O ângulo de 1 rad é aquele em que
  \begin{alternativas}
  \item o arco mede o dobro do raio.
  \item o arco mede a metade do raio.
  \item o arco tem o mesmo comprimento do raio.
  \item a volta completa vale 1 rad.
  \end{alternativas}
\item A volta completa do círculo corresponde a
  \begin{alternativas}
  \item $\pi$ radianos, ou seja, $180^\circ$.
  \item $2\pi$ radianos, ou seja, $360^\circ$.
  \item $90^\circ$, ou seja, $\pi/2$ radianos.
  \item 1 radiano, ou seja, $57^\circ$.
  \end{alternativas}
\item Converta, usando $\pi$ rad $= 180^\circ$ (sem calculadora).
  \begin{alternativas}
  \item $30^\circ$ em radianos
  \item $45^\circ$ em radianos
  \item $180^\circ$ em radianos
  \item $270^\circ$ em radianos
  \item $\pi/3$ em graus
  \item $3\pi/2$ em graus
  \item $2\pi$ em graus
  \end{alternativas}
\item Na corrente alternada, a intensidade $i(t)$
  \begin{alternativas}
  \item mantém o sinal e só muda o módulo.
  \item inverte o sentido a cada meio ciclo.
  \item é nula em qualquer instante $t>0$.
  \item coincide com a corrente de uma pilha.
  \end{alternativas}
\item Converta $90^\circ$ para radianos e diga $\sen 90^\circ$.
\item Complete: $\sen\pi = \ldots$ e $\cos 0 = \ldots$ . O que isso implica
      para $i$ no seno e no cosseno.
\end{questoes}

\nivel{2}{Aplicação}
\begin{questoes}
\item A conversão de $60^\circ$ para radianos dá
  \begin{alternativas}
  \item $\pi/2$, e o seno desse ângulo vale 1.
  \item $\pi/3$, e o seno desse ângulo vale $\sqrt{3}/2$.
  \item $\pi/4$, e o seno desse ângulo vale 1.
  \item $\pi/6$, e o seno desse ângulo vale 0.
  \end{alternativas}
\item Se $i(t) = i_0\sen(\omega t)$ e $\omega t = \pi/2$, então
  \begin{alternativas}
  \item $i = 0$, porque o seno se anula.
  \item $i = -i_0$, pico no sentido contrário.
  \item $i = +i_0$, pico no sentido positivo.
  \item $i = i_0/2$, metade da amplitude.
  \end{alternativas}
\item Uma CA é $i(t) = \dec{2,5}\sen(100\pi\,t)$ (SI). Os valores de $i_0$ e
      $f$ são
  \begin{alternativas}
  \item $i_0 = 100\un{A}$ e $f = \dec{2,5}\un{Hz}$.
  \item $i_0 = \dec{2,5}\un{A}$ e $f = 50\un{Hz}$.
  \item $i_0 = \dec{2,5}\un{A}$ e $f = 100\un{Hz}$.
  \item $i_0 = 50\un{A}$ e $f = \dec{2,5}\un{Hz}$.
  \end{alternativas}
\item Passe $\pi/6$ para graus e calcule $\sen(\pi/6)$.
\item Com $i(t) = \dec{6,0}\cos(\omega t)$ e $\omega t = 0$, a corrente vale
  \begin{alternativas}
  \item 0, pois o cosseno começa nulo.
  \item $-\dec{6,0}\un{A}$, pico negativo.
  \item $+\dec{3,0}\un{A}$, metade do pico.
  \item $+\dec{6,0}\un{A}$, pico positivo.
  \end{alternativas}
\item A rede tem $f = 60\un{Hz}$. Quantos ciclos completos ocorrem em
      $\dec{0,10}\un{s}$? Qual o período $T$?
\end{questoes}

\nivel{3}{Mais trabalhados}
\begin{questoes}
\item Converta $120^\circ$ para radianos. Em seguida, usando
      $\sen 120^\circ = \sen 60^\circ = \sqrt{3}/2$, ache $i$ se
      $i = \dec{4,0}\sen(\omega t)$ e $\omega t = 120^\circ$.
\item Uma CA de amplitude $\dec{8,0}\un{A}$ e $f = 60\un{Hz}$ é
      $i(t) = i_0\sen(2\pi f t)$. Determine: (a) $i(t)$ explícita;
      (b) $i$ em $t = T/4$; (c) o ângulo $\omega t$ em $t = T/4$, em rad e
      em graus.
\tcbbreak
\item Sobre $i(t) = i_0\sen(\omega t)$ e $i(t) = i_0\cos(\omega t)$, assinale
      a correta.
  \begin{alternativas}
  \item são ondas diferentes: uma é CC e a outra é CA.
  \item só o seno descreve CA; o cosseno descreve CC.
  \item descrevem a mesma forma, com origem dos tempos distinta.
  \item o cosseno não pode ser negativo, logo não inverte.
  \end{alternativas}
\item Em $t = \dec{5,0}\times 10^{-3}\un{s}$, $i = \dec{4,0}\sen(100\pi t)$
      vale quanto? Passe $\omega t$ também para graus.
\item Um aluno coloca a calculadora em \dest{graus} e calcula $\sen(\pi/2)$
      como se $\pi/2$ fosse $\dec{1,57}^\circ$. O que sai de errado e como
      consertar?
\item Mostre que $\omega = 2\pi f$ e $f = 1/T$ implicam
      $i(t) = i_0\sen(2\pi t/T)$. Se $T = 1/60\un{s}$, escreva $\omega$
      em rad/s.
\end{questoes}

\gabarito{1c · 2b · 3) (a) $\pi/6$ · (b) $\pi/4$ · (c) $\pi$ · (d) $3\pi/2$ ·
(e) $60^\circ$ · (f) $270^\circ$ · (g) $360^\circ$ · 4b · 5) $\pi/2$;
$\sen = 1$ · 6) 0; 1 · 7b · 8c · 9b · 10) $30^\circ$; $1/2$ · 11d ·
12) 6 ciclos; $T = 1/60\un{s}$ · 13) $2\pi/3$; $i = 2\sqrt{3}\un{A}$ ·
14) $i = \dec{8,0}\sen(120\pi t)$; $+\dec{8,0}\un{A}$; $\pi/2 = 90^\circ$ ·
15c · 16) $\omega t = \pi/2 = 90^\circ \Rightarrow i = \dec{4,0}\un{A}$ ·
17) $\pi/2$ é radiano ($90^\circ$); em graus daria seno quase nulo ·
18) $\omega = 120\pi$ rad/s.}
\end{exercicios}

\end{document}
